\begin{minipage}{0.68\linewidth}
On fait créer une onde circulaire dans la cuve à ondes, on ajuste la fréquence
de l'onde circulaire à différentes valeurs, et à chaque fois on éclaire la
surface de l'eau avec un stroboscope réglé à la même fréquence de l'onde, on
observe que tous les points de la surface de l'eau apparaissent immobiles, puis
on mesure la longueur d'onde correspondante.
    
\begin{enumerate}
    \item Compléter le tableau suivant~:
          \begin{table}[H]
              \centering
              \setlength{\tabcolsep}{18pt}
              \renewcommand{\arraystretch}{1.5}
              \begin{tabularx}{\textwidth}{|p{1.5cm}|X|X|X|X|}
                  \hline $f~(\unit{\hertz})$ & 20 & 25 & 30 & 35 \\
                  \hline $\lambda~(\unit{\m})$ & 1 & 0,9 & 0,8 & 0,7 \\
                  \hline $v~(\unit{\v})$ &&&& \\
                  \hline
              \end{tabularx}
          \end{table}
    \item On dit qu'un milieu est dispersif si la vitesse de propagation d'une
          onde dans ce milieu dépend de sa fréquence. L'eau est-elle un milieu
          dispersif~?
\end{enumerate}
\end{minipage}
\hfill
\begin{minipage}{0.30\linewidth}
    % \fulshright
    \begin{figure}[H]
    \includegraphics[width=0.8\linewidth]{cuve_a_onde}
    \caption{Cuve à onde}
\end{figure}
\end{minipage}
